\section{下一步工作的方向建议}

本章把建议分成四条轨道：验证轨道（最短路径上的核验）、数学推进轨道（问题本身往哪走）、方法回灌轨道（这套装置能否搬家）、工程流程轨道（共同体怎么接住批量产出）。每条建议都锚定手稿明确留下的空白或本文核实过的缺口，并附可行性、成本与收益的量级判断；面向研究生与青年学者的可进入点集中放在最后。

\subsection{轨道一：验证}

\begin{dirbox}[验证轨道]
\textbf{1）模块化审读。}按第 9.4 节的 M1--M8 拆分，先把 M1 与链条根（三维引理段 + Guth--Wang--Zahl 原文）做掉。可行性：高；成本：约 4--6 人月的专家工时（按每模块 20--35 页的密度估，\cfM）；收益：给出“论证是否值得继续投入”的第一道闸门。\\
\textbf{2）局部 Lean 形式化。}优先 M1 与三维引理段；不追求全覆盖，追求“依赖链根部可信”。参照人类侧先例（三维挂谷形式化\upcite{numina2026}）；同时配套做翻译保真度校验，避免自动形式化的语义漂移\upcite{bastounis2026}。可行性：中；成本：高于人工审读、低于全篇形式化；收益：一劳永逸的机械核验层。\\
\textbf{3）离散模型压力测试。}对三腿切换（$2k=1$）与次数恢复（圆锥障碍）构造有限域类比或小规模数值实验，检验机制在退化情形下的稳定性。可行性：高（适合作为学位论文章节）；成本：低；收益：为“机制可信度”提供独立于文本的证据。
\end{dirbox}

\subsection{轨道二：数学推进}

手稿留下的空白是密集的，按“离手稿越近越优先”排序：

\begin{dirbox}[数学推进轨道]
\textbf{1）四维极大函数猜想。}手稿只做了 Hausdorff 维数；二维、三维的极大版本均已解决\upcite{cordoba1977,openai3d2026}，四维极大是同一语法下的下一个未结命题。可行性：低（难度即刻上升一档）；成本：高；收益：若成，将与四维集合定理合成完整图景。\\
\textbf{2）$n\ge 5$ 的框架推广。}手稿目前只给投影下界（$\dim_{\mathrm{H}}\ge 4$）\upcite{openai4d2026}。一个自然的猜想式路线：把“二次图卡”升级为“低次代数曲面图卡”，并把窄度参数 $\ell$ 的区间划分改写为按“拟合次数 vs 时间尺度”的分层——四维里 $\ell=0$ 对应的圆锥障碍，在高维将替换为更广的簇障碍。可行性：未知；成本：极高。\\
\textbf{3）概念蒸馏与再表述。}把地平线、两色匹配、水平模型写成独立的、可讲授的材料（讨论班讲义、教程），检验这些概念是否独立成立、是否可以被改造后用于同类问题（如投影理论、Radon 型变换）。可行性：高；成本：低；收益：直接服务“理解账”，且这类工作天然适合作为综述与教科书的章节。\\
\textbf{4）定量联动。}三维篇已给出极大估计的 $\delta^{-\varepsilon}$ 版本，与 Fourier 限制、Bochner--Riesz、局部光滑的联动在三维可以立即着手；球面限制旁路篇（family 077）与挂谷主线的关系值得专文对照\upcite{openaielliptic2026,openai3d2026,gwz2026}。
\end{dirbox}

\subsection{轨道三：方法回灌}

\begin{dirbox}[方法回灌轨道]
\textbf{1）地平线框架的移植。}“局部拟合精度 $\leftrightarrow$ 时间代价”的记账方式不依赖四维几何本身；把同一装置移植到：限制型算子的波包聚集问题、Furstenberg 型投影问题、以及 Kakeya--Nikodym 的极值分析，是低风险高观测价值的试探。可行性：中；成本：中。\\
\textbf{2）熵账本与平面刷的融合。}门内的平面性/颗粒性机制（平面刷）与借入的熵/耦合机制目前是两套语言；若能用一套统一表述覆盖两者，将显著降低后续论文的阅读成本。可行性：中；成本：中；收益：为下一代工作提供公共语言。\\
\textbf{3）把“打包—复用”变成常规。}Guth--Wang--Zahl 流线化 → OpenAI 三维加权化 → 四维复用，这条打包链展示了可组合性红利；建议将“引理接口的稳定性”当作独立指标来维护（例如在形式化库里固定接口），这也回应了第 9 章的验证瓶颈。可行性：中高；成本：中。
\end{dirbox}

\subsection{轨道四：工程流程}

\begin{dirbox}[工程流程轨道]
\textbf{1）分诊协议。}对批量产出建立公开的分诊表：按“重要性 × 可验证性 × 人力预算”排序，先做依赖链根部与高复用潜力的结果。本综述的附录 B 即是一份微型样例。\\
\textbf{2）形式化覆盖率作为发布契约。}仓库已按版本管理修订（撤稿、补丁、形式化增量）\upcite{openairepo2026}；社区侧的合理诉求是把“顶层结果形式化覆盖率”与“独立审读状态”纳入发布标准——AHM 声明要求的“停止合作”与 OpenAI 主张的“继续对话”之间，至少存在一个技术性共识区：验证基础设施的共建\upcite{ahm2026,openairelease2026,tao2026}。\\
\textbf{3）评价指标的补齐。}当前对这类发布，唯一可公开核对的指标是形式化覆盖率（300/719 $\approx$ 42\%）；真正决定长期影响的“独立审读周期”与“概念复用率”尚无任何制度性跟踪。建议以 074 族为试点建立这两项记录。
\end{dirbox}

\subsection{面向研究生的可进入点}

\begin{table}[t]
\centering
\caption{具体可进入的小课题与验收标准}
\label{tab:entrypoints}
\footnotesize
\begin{tabularx}{\textwidth}{@{}L{0.6cm}L{4.6cm}Y@{}}
\toprule
\rowcolor{lightbg}\textbf{编号} & \textbf{题目} & \textbf{验收标准} \\
\midrule
E1 & 三维篇加权板条引理的独立复演（含与 Guth--Wang--Zahl 并集估计的接口） & 交付一份自足推导（10--20 页），专家复核无缺环；并记录所有用到的假设 \\
E2 & 四维篇归约（M1）假设审计 & 逐步清单：每一步的假设、结论、隐含常数依赖；确认无测度性假设混入 \\
E3 & 三腿切换的离散模型实验 & 可复现代码 + 报告：在 $\mathbb{F}_p$ 类比或数值模型中观察/反驳机制；记录失败情形 \\
E4 & Proposition 6.8 的交互式形式化 & 可编译的局部形式化脚本（允许带未证洞），配翻译对照表 \\
E5 & 推论章转移定理条件核对（第 11.4--11.5 节） & 条件对照表：每条转移定理的假设 vs 手稿使用处，逐条标绿/标红 \\
\bottomrule
\end{tabularx}
\end{table}

E1 与 E2 是“小投入快反馈”的一对：不要求核验者同意手稿的全部机制，只要求把链条根部的推导独立走一遍——数学上干净，方法上保守，产出立刻可被共同体使用。E3 适合有计算背景的学生；E4 适合形式化方向；E5 适合偏好文献与结构分析的工作。

\begin{keybox}[可记忆主句]
这份手稿给学界留下的最有价值的东西，可能不是那个待验的四维结论，而是一张已经画好的攻击地图：谁先按图索骥核验链条根部，谁就掌握了后续所有工作的闸门。
\end{keybox}